\documentclass[12pt]{article}
\usepackage[utf8]{inputenc}
\usepackage{graphicx}
\usepackage{enumerate}
\newcommand{\pt}[1]{\textbf{(#1 pts)}}
\begin{document}
\includegraphics[scale=.25]{LogoMate.eps}

\vspace{-3ex}
$\phantom{x}$\hfill MAT003 Cálculo II\\
$\phantom{x}$\hfill Laboratorio 1, II Ciclo 2025\\[1ex]
Facultad de Ciencias Exactas y Naturales\hfill Duración: 2 horas\\
Lunes  1 de setiembre de 2025\\
Puntaje total: 4 puntos
\vspace{-1ex}
\hrulefill
\hfill v1 - SOLUCIONES

\section*{Pregunta 1 \pt{4}}
Calcule las siguientes medidas para la variable Peso: Q2 y media.

\paragraph{Solución y retroalimentación:} El cuartil 2 (Q2) de la variable Peso corresponde al percentil 50. Para determinarlo se consideran las 100 observaciones ordenadas de menor a mayor y se obtiene el valor asociado con el 50 \% de la distribución. El resultado es 69,40 kg. Esto indica que aproximadamente el 50 \% de las observaciones presentan valores de peso menores o iguales a 69,40 kg.

La media aritmética de la variable Peso se calcula sumando sus 100 observaciones y dividiendo el resultado entre 100. La suma es 6966,50; por tanto, la media es 69,67 kg. Esto representa el promedio de peso en las observaciones.
\end{document}